\section{Refinement Reflection: \corelan}
\label{sec:formalism}
\label{sec:types-reflection}
\label{sec:theory}

We formalize refinement reflection in two steps.
%
First, we develop a core calculus \corelan with an
\emph{undecidable} type system based on denotational
semantics. We show how the soundness of the type system
allows us to \emph{prove theorems} using \corelan.
%
Next, in \S~\ref{sec:theory:algorithmic} we define
a language \smtlan that soundly approximates \corelan
while enabling decidable SMT-based type checking.


\subsection{Syntax}

\begin{figure}[t]
\begin{minipage}[b]{0.47\textwidth}
%\newcommand{\ra}[1]{\renewcommand{\arraystretch}{#1}}

%\begin{figure}[t!]
%\vspace{-5mm}
%\centering
$$
\begin{array}{rrcl}
\emphbf{Operators}
  & \odot
  & ::= & = \spmid  <
\\[0.03in]

\emphbf{Constants}
  & c
  & ::=
  & \land \spmid \lnot \spmid \odot \spmid +,-,\dots  \\
  && \spmid & \etrue\spmid \efalse \spmid 0, \pm 1, \dots
\\[0.03in]

\emphbf{Values}
  & w & ::=&  c
             \spmid \efun{x}{\typ}{e} \spmid D\ \overline{w}
\\[0.03in]

\emphbf{Expressions}
  & e & ::=    & w \spmid x \spmid \eapp{e}{e}      \\
  &   & \spmid & \ecase{x}{e}{\dc}{\overline{x}}{e}
  %% &   & \spmid & \eletb{x}{\gtyp}{e}{e}
\\[0.03in]

\emphbf{Binders}
  & \bd & ::= & e \spmid \eletb{x}{\gtyp}{\bd}{\bd}
\\[0.03in]

\emphbf{Program}
  & \prog & ::= & \bd \spmid \erefb{x}{\gtyp}{e}{\prog}
\\[0.03in]

% \emphbf{Logical Labels}
%  & L & ::= & \la \spmid \lm
%\\[0.03in]

\emphbf{Basic Types}
  & \btyp
  & ::=
  & \tint \spmid \tbool \spmid T
\\[0.03in]

\emphbf{Ref. Types}
  & \typ
  & ::=   & \tref{v}{\btyp^{[\tlabel]}}{\reft} \spmid \tfun{x}{\typ}{\typ}
\\[0.05in]
\end{array}
$$
\caption{\textbf{Syntax of \corelan}}
\label{fig:syntax}

%\vspace{-2mm}
% \end{figure}

%%\begin{figure}[t!]
%%\centering
%%
%%\emphbf{Contexts}\hfill{{$\fbox{\textit{C}}$}}
%%$$
%%\begin{array}{rcl}
%%  C & ::=    & \bullet \spmid C\ e \spmid c\ C \spmid D\ \overline{e}\ C\ \overline{e}\\
%%    & \spmid & \ecase{y}{C}{\dc_i}{\overline{x}}{e_i}
%%  \\[0.03in]
%%\end{array}
%%$$
%%
%%\emphbf{Reductions}\hfill{{$\fbox{\goesto{\prog}{\prog'}}$}}
%%\begin{align*}
%%C[\prog]
%%  & \hookrightarrow
%%   C[\prog'],\quad \text{if}\ \goesto{\prog}{\prog'}
%%  \\
%%{c\ v}
%%  & \hookrightarrow
%%   {\ceval{c}{v}}
%%  \\
%%{({\efun{x}{\typ}{e})}\ {e'}}
%%  & \hookrightarrow
%%   {\SUBST{e}{x}{e'}}
%%  \\
%%{\ecase{y}{\dc_j\ \overline{e}}{\dc_i}{\overline{x_i}}{e_i}}
%%  &\hookrightarrow
%%  {\SUBST{\SUBST{e_j}{y}{\dc_j\ \overline{e}}}{\overline{x_i}}{\overline{e}}}
%%\\
%%{\erefb{x}{\gtyp}{e}{\prog}}
%%  & \hookrightarrow
%%  {\SUBST{\prog}{x}{\efix{(\efun{x}{\gtyp}{e})}}}
%%\\
%%{\eletb{x}{\gtyp}{\bd_x}{\bd}}
%%  & \hookrightarrow
%%{\SUBST{\bd}{x}{\efix{(\efun{x}{\gtyp}{\bd_x})}}}
%%\\
%%{\efix{\prog}}
%%  & \hookrightarrow
%%  {\prog\ (\efix{\prog})}
%%\end{align*}
%%\caption{\textbf{Operational Semantics of \corelan}}
%%\label{fig:semantics}
%%\end{figure}

\end{minipage}
\hspace{0.2in}
\begin{minipage}[b]{0.47\textwidth}
%% \begin{figure}[t!]
%% \vspace{-5mm}
%% \centering
$$
\begin{array}{rrcl}
\emphbf{Predicates}
  & \pred & ::= &
    \pred \binop \pred \spmid
    \unop \pred \\
  && \spmid & n \spmid b \spmid x \spmid \dc \spmid  x\ \overline{\pred}\\
%%  && \spmid & \forall \overline{\tbind{x}{\sort}}. \pred
  && \spmid & \eif{\pred}{\pred}{\pred}
\\[0.03in]

\emphbf{Integers}
  & n
  & ::= & 0, -1, 1, \dots
\\[0.03in]

\emphbf{Booleans}
  & b
  & ::= & \etrue \spmid \efalse
\\[0.03in]

\emphbf{Binary Ops}
  & \binop
  & ::= & = \spmid < \spmid \land \spmid + \spmid - \spmid \dots
\\[0.03in]

\emphbf{Unary Ops}
  & \unop
  & ::= & \lnot \spmid \dots
\\[0.03in]

\emphbf{Sort Args}
  & \sort_a
  & ::= & \tint \spmid \tbool \spmid \tuniv
         \spmid \tsmtfun{\sort_a}{\sort_a}
\\[0.03in]
\emphbf{Sorts}
  & \sort
  & ::=  & \sort_a \spmid \sort_a \rightarrow \sort
\end{array}
$$
\caption{\textbf{Syntax of \smtlan}}
\label{fig:smtsyntax}
% \vspace{-2mm}
% \end{figure}

\end{minipage}
\end{figure}

%
Figure~\ref{fig:syntax} summarizes the syntax of \corelan,
which is essentially the calculus \undeclang~\cite{Vazou14}
with explicit recursion and a special $\erefname$ binding
to denote terms that are reflected into the refinement logic.
%
The elements of \corelan are layered into
primitive constants, values, expressions, binders
and programs.

\mypara{Constants}
The primitive constants of \corelan
include primitive relations $\op$,
here, the set $\{ =, <\}$.
%
Moreover, they include the
primitive booleans $\etrue$, $\efalse$,
integers $\mathtt{-1}, \mathtt{0}$, $\mathtt{1}$, \etc,
and logical operators $\mathtt{\land}$, $\mathtt{\lor}$, $\mathtt{\lnot}$, \etc.

\mypara{Data Constructors}
%
Data constructors are special constants.
% Each data type has an equality predicate $\haseq{T}$
% that is true only if values of type $T$ can be finitely compared.
For example, the data type \tintlist, which represents
finite lists of integers, has two data constructors: $\dnull$ (nil)
and $\dcons$ (cons).

\mypara{Values \& Expressions}
%
The values of \corelan include
constants, $\lambda$-abstractions
$\efun{x}{\typ}{e}$, and fully
applied data constructors $D$
that wrap values.
%
The expressions of \corelan
include values, variables $x$,
applications $\eapp{e}{e}$, and
$\mathtt{case}$ expressions.

\mypara{Binders \& Programs}
%
A \emph{binder} $\bd$ is a series of possibly recursive
$\mathtt{let}$ definitions, followed by an expression.
%
A \emph{program} \prog is a series of $\erefname$
definitions, each of which names a function
that is reflected into the refinement
logic, followed by a binder.
%
The stratification of programs via binders
is required so that arbitrary recursive definitions
are allowed in the program but cannot be inserted into the logic
via refinements or reflection.
%
(We \emph{can} allow non-recursive $\mathtt{let}$
binders in expressions $e$, but omit them for simplicity.)

\subsection{Operational Semantics}
We define $\hookrightarrow$ to be the small step, call-by-name
$\beta$-reduction semantics for \corelan.
%
We evaluate reflected terms %$\erefb{x}{\gtyp}{e}{\prog}$
as recursive $\mathtt{let}$ bindings, with extra termination-check
constraints imposed by the type system:
%
$$
\erefb{x}{\gtyp}{e}{\prog}
\hookrightarrow
\eletb{x}{\gtyp}{e}{\prog}
$$
%
We define $\evalsto{}{}$ to be the reflexive,
transitive closure of $\evals{}{}$.
%
Moreover, we define $\betaeq{}{}$ to be the reflexive,
symmetric, and transitive closure of $\evals{}{}$.

\mypara{Constants} Application of a constant requires the
argument be reduced to a value; in a single step, the
expression is reduced to the output of the primitive
constant operation, \ie ${c\ v} \hookrightarrow \ceval{c}{v}$.
%
For example, consider $=$, the primitive equality
operator on integers.
%
We have $\ceval{=}{n} \defeq =_n$
where $\ceval{=_n}{m}$ equals \etrue
iff $m$ is the same as $n$.
%

\mypara{Equality}
We assume that the equality operator
is defined \emph{for all} values,
and, for functions, is defined as
extensional equality.
%
That is, for all
$f$ and $f'$,
$\evals{(f = f')}{\etrue}$
   $\mbox{iff}$
  $\forall v.\ \betaeq{f\ v}{f'\ v}$.
%
We assume source \emph{terms} only contain implementable equalities
over non-function types; while function extensional equality only
appears in \emph{refinements}. % and is approximated by the underlying logic.

\subsection{Types}

\corelan types include basic types, which are \emph{refined} with predicates,
and dependent function types.
%
\emph{Basic types} \btyp comprise integers, booleans, and a family of data-types
$T$ (representing lists, trees \etc).
%
For example, the data type \tintlist represents lists of integers.
%
We refine basic types with predicates (boolean-valued expressions \refa) to obtain
\emph{basic refinement types} $\tref{v}{\btyp}{\refa}$.
%
%
We use \tlabel to mark provably terminating
computations and use
refinements to ensure that if
${{e}\text{:}{\tref{v}{\btyp^\tlabel}{\refa'}}}$,
then $e$ terminates.
%
%OK \NV{NEW}
As discussed by~\citet{Vazou14} termination labels
can be checked using refinement types and are used
to ensure that refinements cannot diverge as required
for soundness of type checking under lazy evaluation.
%OK \NV{END NEW}
%
Finally, we have dependent \emph{function types} $\tfun{x}{\typ_x}{\typ}$
where the input $x$ has the type $\typ_x$ and the output $\typ$ may
refer to the input binder $x$.
%
We write $\btyp$ to abbreviate $\tref{v}{\btyp}{\etrue}$,
and \tfunbasic{\typ_x}{\typ} to abbreviate \tfun{x}{\typ_x}{\typ} if
$x$ does not appear in $\typ$.



\mypara{Denotations}
%
Each type $\typ$ \emph{denotes} a set of expressions $\interp{\typ}$,
that is defined via the operational semantics~\cite{Knowles10}.
%
Let $\shape{\typ}$ be the type we get if we erase all refinements
from $\typ$ and $\bhastype{}{e}{\shape{\typ}}$ be the
standard typing relation for the typed lambda calculus.
%
Then, we define the denotation of types as:

\begin{align*}
\interp{\tref{x}{\btyp}{r}} \defeq &\
  \{ e \mid  \bhastype{}{e}{\btyp}
           , \mbox{ if } \evalsto{e}{w} \mbox{ then } \evalsto{r\subst{x}{w}}{\etrue}
  \}\\
\interp{\tref{x}{\btyp^\tlabel}{r}} \defeq &\
  \interp{\tref{x}{\btyp}{r}} \cap \{ e \mid  \evalsto{e}{w}\}\\
\interp{\tfun{x}{\typ_x}{\typ}} \defeq &\
  \{ e \mid \bhastype{}{e}{\shape{\tfunbasic{\typ_x}{\typ}}}
           , \forall e_x \in \interp{\typ_x}.\ (\eapp{e}{e_x}) \in \interp{\typ\subst{x}{e_x}}
  \}
\end{align*}


\mypara{Constants}
For each constant $c$ we define its type \constty{c}
such that $c \in \interp{\constty{c}}$.
%
For example,
%
$$
\begin{array}{lcl}
\constty{3} &\doteq& \tref{v}{\tint^\tlabel}{v = 3}\\
\constty{+} &\doteq& \tfun{\ttx}{\tint^\tlabel}{\tfun{\tty}{\tint^\tlabel}{\tref{v}{\tint^\tlabel}{v = x + y}}}\\
\constty{\leq} &\doteq& \tfun{\ttx}{\tint^\tlabel}{\tfun{\tty}{\tint^\tlabel}{\tref{v}{\tbool^\tlabel}{v \Leftrightarrow x \leq y}}}\\
\end{array}
$$


\subsection{Refinement Reflection}\label{subsec:logicalannotations}
%
The key idea in our work is to
\emph{strengthen} the output type of functions
with a refinement that \emph{reflects} the
definition of the function in the logic.
%
We do this by treating each
%
$\erefname$-binder
%
(${\erefb{f}{\gtyp}{e}{\prog}}$)
%
as a $\eletname$-binder
%
(${\eletb{f}{\exacttype{\gtyp}{e}}{e}{\prog}}$)
%
during type checking (rule $\rtreflect$ in Figure~\ref{fig:typing}).

\mypara{Reflection}
%
We write \exacttype{\typ}{e} for the \emph{reflection}
of the term $e$ into the type $\typ$,  defined by strengthening
\typ as:
%
$$
\begin{array}{lcl}
\exacttype{\tref{v}{\btyp{^\tlabel}}{r}}{e}
  & \defeq
  & \tref{v}{\btyp{^\tlabel}}{r \land v = e}\\
\exacttype{\tfun{x}{\typ_x}{\typ}}{\efun{x}{}{e}}
  & \defeq
  & \tfun{x}{\typ_x}{\exacttype{\typ}{e}}
\end{array}
$$
%
As an example, recall from \S~\ref{sec:overview}
that the @reflect fib@ strengthens the type of
@fib@ with the refinement @fibP@.

% OK \NV{NEW}

\mypara{Termination Checking}
We defined \exacttype{\cdot}{\cdot} to be a \textit{partial}
function that only reflects provably terminating expressions,
\ie expressions whose result type is marked with \tlabel.
%
If a non-provably terminating function is reflected
in an \corelan expression then type checking will
fail (with a reflection type error in the implementation).
%
This restriction is crucial for soundness,
as diverging expressions can lead to inconsistencies.
%
For example, reflecting the diverging @f x = 1 + f x@
into the logic leads to an inconsistent system that
is able to prove @0 = 1@.
%%This choice complies with both the facts that
%%all refinement expressions are checked to be terminating
%%by the type system (by the rule \rwbase) and, as explained in \S~\ref{sec:algorithmic},
%%all reflected functions are encoded as total uninterpreted functions
%%in the SMT logic of \smtlan.

\mypara{Automatic Reflection}
Reflection of \corelan expressions into the refinements happens
automatically by the type system, not manually by the user.
%
The user simply annotates a function $f$ as $\annotReflect\ f$.
%
Then, the rule $\rtreflect$ in Figure~\ref{fig:typing}
is used to type check the reflected function by strengthening
the $f$'s result via \exacttype{\cdot}{\cdot}.
%
Finally, the rule \rtlet is used to check that the
automatically strengthened type of $f$ satisfies $f$'s implementation.

% OK \NV{END NEW}

\mypara{Consequences for Verification}
%
% Reflection has two consequences for verification.
Reflection has a major consequence for verification:
%
%First, the reflected refinement is \emph{not trusted};
%it is itself verified (as a valid output type)
%during type checking.
%%
%Second,
instead of being tethered to quantifier
instantiation heuristics or having to
program ``triggers'' as in
Dafny~\citep{dafny} or \fstar~\citep{fstar},
%
the programmer can predictably ``unfold'' the
definition of the function during a proof simply
by ``calling'' the function, which, as discussed
in~\S~\ref{sec:evaluation}, we have found to be
a very natural way of structuring proofs.




\subsection{Typing Rules}
\begin{figure}[!t]
\begin{minipage}[t]{0.47\textwidth}
  \emphbf{Typing}\hfill{\fbox{\hastype{\env}{\prog}{\typ}}}\\
$$
\inference{
	\tbind{x}{\gtyp}\in\env
}{
	\hastype{\env}{x}{\gtyp}
}[\rtvar]
\qquad
\inference{
}{
	\hastype{\env}{c}{\constty{c}}
}[\rtconst]
$$

$$
\inference{
	\hastype{\env}{\prog}{\typ'} &
	\issubtype{\env}{\typ'}{\typ}
}{
	\hastype{\env}{\prog}{\typ}
}[\rtsub]
$$
$$
\inference{
    % \haseq{\btyp} &
	\hastype{\env}{e}{\tref{v}{\btyp}{\reft_r}}
}{
	\hastype{\env}{e}{\tref{v}{\btyp}{\reft_r\land v = e}}
}[\rtexact]
$$
$$
\inference{
	\hastype{\env, \tbind{x}{\typ_x}}{e}{\typ}
}{
	\hastype{\env}{\efun{x}{\typ}{e}}{\tfun{x}{\typ_x}{\typ}}
}[\rtfun]
$$

$$
\inference{
	\hastype{\env}{e_1}{(\tfun{x}{\typ_x}{\typ})} &&
	\hastype{\env}{e_2}{\typ_x}
}{
	\hastype{\env}{e_1\ e_2}{\typ}
}[\rtapp]
$$

$$
\inference
	{\hastype{\env, \tbind{x}{\gtyp_x}}{\bd_x}{\gtyp_x} &
	 \iswellformed{\env, \tbind{x}{\gtyp_x}}{\typ_x} \\
	 \hastype{\env, \tbind{x}{\gtyp_x}}{\bd}{\gtyp} &
	 \iswellformed{\env}{\typ}}
	{\hastype{\env}{\eletb{x}{\gtyp_x}{\bd_x}{\bd}}{\typ}}
	[\rtlet]
$$

$$
\inference
	{\hastype{\env}
	 				 {\eletb{f}{\exacttype{\gtyp_f}{e}}{e}{\prog}}
					 {\typ}
	}
	{\hastype{\env}
					 {\erefb{f}{\gtyp_f}{e}{\prog}}
					 {\typ}
	}[\rtreflect]
$$

$$
\inference{
	\hastype{\env}{e}{\tref{v}{T}{e_r}} & \iswellformed{\env}{\typ} \\
	& \forall i. \constty{\dc_i} = \tfunbasic{\overline{\tbind{y_j}{\typ_j}}}{\tref{v}{T}{e_{r_i}}} \\
	& \hastype{\env, \overline{\tbind{y_j}{\typ_j}}, \tbind{x}{\tref{v}{T}{e_r \land e_{r_i}}} }{e_i}{\typ}
}{
	\hastype{\env}{\ecase{x}{e}{\dc_i}{\overline{y_i}}{e_i}}{\typ}
}[\rtcase]
$$

\end{minipage}
\hspace{0.2in}
\begin{minipage}[t]{0.47\textwidth}
	\emphbf{Well Formedness}\hfill{\fbox{\iswellformed{\env}{\typ}}}\\

$$
\inference{
  \hastype{\env,\tbind{v}{\btyp}}{\refa}{\tbool^{\tlabel}}
}{
	\iswellformed{\env}{\tref{v}{\btyp}{\refa}}
}[\rwbase]
$$
$$
\inference{
	\iswellformed{\env}{\typ_x} &
	\iswellformed{\env,\tbind{x}{\typ_x}}{\typ}
}{
	\iswellformed{\env}{\tfun{x}{\typ_x}{\typ}}
}[\rwfun]
$$

	\emphbf{Subtyping}\hfill{\fbox{\issubtype{\env}{\typ_1}{\typ_2}}}\\

$$
\inference{
        \forall \sub\in\interp{\env}.
        \interp{\applysub{\sub}{\tref{v}{\btyp}{\refa_1}}}
        \subseteq
        \interp{\applysub{\sub}{\tref{v}{\btyp}{\refa_2}}}
}{
        \issubtype{\env}{\tref{v}{\btyp}{\refa_1}}{\tref{v}{\btyp}{\refa_2}}
}[\rsubbasecore]
$$

$$
\inference{
	\issubtype{\env}{\typ_x'}{\typ_x} &
	\issubtype{\env,\tbind{x}{\typ_x'}}{\typ}{\typ'}
}{
	\issubtype{\env}{\tfun{x}{\typ_x}{\typ}}{\tfun{x}{\typ_x'}{\typ'}}
}[\rsubfun]
$$\\[0.4in]

\emphbf{Algorithmic-Subtyping}\hfill{\fbox{\aissubtype{\env}{\typ_1}{\typ_2}}}\\

$$
\inference{
%\env' \defeq \env,\tbind{v}{\{\btyp^\tlabel | \refa\}} \\
\env' \defeq \env,\tbind{v}{\tref{v}{\btyp^\tlabel}{\refa}} \\
\tologicshort{\env'}{\refa'}{\tbool}{\pred'}{}{}{} &
\smtvalid{\vcond{\env'}{\pred'}}
%%	\forall \sub\in\interp{\env}.
%%	\interp{\applysub{\sub}{\tref{v}{\btyp}{\refa_1}}}
%%	\subseteq
%%	\interp{\applysub{\sub}{\tref{v}{\btyp}{\refa_2}}}
}{
	\aissubtype{\env}{\tref{v}{\btyp}{\refa}}{\tref{v}{\btyp}{\refa'}}
}[\rsubbasesmt]
$$

\end{minipage}
\caption{Typing of \corelan and algorithmic subtyping of \smtlan}
\label{fig:typing}
\end{figure}

%
Next, we present the type-checking rules of \corelan.

\mypara{Environments and Closing Substitutions}
A \emph{type environment} $\env$ is a sequence of type bindings
$\tbind{x_1}{\typ_1},\ldots,\tbind{x_n}{\typ_n}$. An environment
denotes a set of \emph{closing substitutions} $\sto$ which are
sequences of expression bindings:
$\gbind{x_1}{e_1}, \ldots,$ $\gbind{x_n}{e_n}$ such that:
$$
\interp{\env} \defeq  \{\sto \mid \forall \tbind{x}{\typ} \in \Env.
                                    \sto(x) \in \interp{\applysub{\sto}{\typ}} \}
$$



\mypara{Typing}
A judgment \hastype{\env}{\prog}{\typ} states that
the program $\prog$ has the type $\typ$ in
the environment $\env$.
That is, when the free variables in $\prog$ are
bound to expressions described by $\env$, the
program $\prog$ will evaluate to a value
described by $\typ$.

\mypara{Rules}
%
All but two of the rules are standard~\cite{Knowles10,Vazou14}.
%
First, rule \rtreflect is used to strengthen the type of each
reflected binder with its definition, as described previously
in \S~\ref{subsec:logicalannotations}.
%
Second, rule \rtexact strengthens the expression with
a singleton type equating the value and the expression
(\ie reflecting the expression in the type).
%
This is a generalization of the ``selfification'' rules
from \cite{Ou2004,Knowles10} and is required to
equate the reflected functions with their definitions.
%
For example, the application @fib 1@ is typed as
${\tref{v}{\tint^\tlabel}{\fibdef\ v\ 1 \wedge v = \fib\ 1}}$ where
the first conjunct comes from the (reflection-strengthened)
output refinement of @fib@~\S~\ref{sec:overview} and
the second comes from rule~\rtexact.
%
%%Finally, rule \rtfix is used to type the intermediate
%%$\texttt{fix}$ expressions that appear, not in the
%%surface language but as intermediate terms in the
%%operational semantics.

\mypara{Well-formedness}
%
A judgment \iswellformed{\env}{\typ} states that
the refinement type $\typ$ is well-formed in
the environment $\env$.
%
Following~\citet{Vazou14}, $\typ$ is well-formed if all
the refinements in $\typ$ are $\tbool$-typed,
provably terminating expressions in $\env$.

%%\mypara{Termination}
%%%
%%Under arbitrary beta-reduction semantics
%%(which includes lazy evaluation), soundness
%%of refinement type checking requires checking
%%termination, for two reasons:
%%%
%%(1)~to ensure that refinements cannot diverge, and
%%(2)~to account for the environment during subtyping~\citep{Vazou14}.
%%%
%%We use \tlabel to mark provably terminating
%%computations, and extend the rules to use
%%refinements to ensure that if
%%${\ahastype{\env}{e}{\tref{v}{\btyp^\tlabel}{r}}}$,
%%then $e$ terminates~\citep{Vazou14}.

\mypara{Subtyping}
%
A judgment \issubtype{\env}{\typ_1}{\typ_2} states
that the type $\typ_1$ is a subtype of %the type
$\typ_2$ in the environment $\env$.
%
%%Informally, $\typ_1$ is a subtype of $\typ_2$ if
%%the refinement of $\typ_1$ \emph{implies}
%%the refinement of $\typ_2$
%%under the assumptions described by $\env$.
%
Informally, $\typ_1$ is a subtype of $\typ_2$ if, when
the free variables of $\typ_1$ and $\typ_2$
are bound to expressions described by $\env$,
the denotation of $\typ_1$
is \emph{contained in} the denotation of $\typ_2$.
%
Subtyping of basic types reduces to denotational
containment checking, shown in rule \rsubbasecore.
%
That is, $\typ_1$ is a subtype of $\typ_2$ under $\env$
if for any closing substitution $\sto$ in the denotation
of $\env$,
$\interp{\applysub{\sto}{\typ_1}}$ is contained in
$\interp{\applysub{\sto}{\typ_2}}$.

\mypara{Soundness}
%
Following \undeclang~\citep{Vazou14},
in ~\citet{appendix}
we prove that evaluation preserves typing
and typing implies denotational
inclusion.
%
\begin{theorem}{[Soundness of \corelan]}\label{thm:safety}
\begin{itemize}
\item\textbf{Denotations}
If $\hastype{\env}{\prog}{\typ}$ then
$\forall \sto\in \interp{\env}. \applysub{\sto}{\prog} \in \interp{\applysub{\sto}{\typ}}$.
\item\textbf{Preservation}
If \hastype{\emptyset}{\prog}{\typ}
       and $\evalsto{\prog}{w}$ then $\hastype{\emptyset}{w}{\typ}$.
\end{itemize}
\end{theorem}

Theorem~\ref{thm:safety} lets us interpret well
typed programs as proofs of propositions.
%
For example, in \S~\ref{sec:overview} we verified
that the term \fibincrname proves
$${\tfun{n}{\tnat}{\ttref{\fib{n} \leq \fib{(n+1)}}}}$$
%
Via soundness of \corelan, we get that for each valid input @n@,
the result refinement is valid.
$$
\forall n. \evalsto{0 \leq n}{\etrue} \Rightarrow \evalsto{\texttt{fib}\ n \leq \texttt{fib}\ {(n+1)}}{\etrue}
$$
